\paragraph{Calibration is not informativeness.} A predictive distribution can be calibrated, sharp and correlated with its own errors without using its input. In EEG-to-fMRI synthesis this is the default, because anatomical variability decides both where errors occur and where any model is uncertain. All eleven methods learn this structure, and so does the EEG-free template, which none of them beats.

\paragraph{Heteroscedastic training can blind a synthesizer to its input.} A heteroscedastic likelihood can explain a residual by widening its variance instead of moving its mean, and with a weakly informative input and a high learning rate it does \citep{seitzer2022pitfalls}. Here this blinded a backbone that otherwise recovers an injected EEG signal, while its calibration stayed respectable; only a positive control reveals it.

\paragraph{What the EEG predicts physiology.} At the voxel level, models that can see an EEG-coupled signal find only the shared task timing. At the region level a lagged linear model predicts motor BOLD as convincingly as published results, but heart rate predicts it as well, the EEG adds nothing to heart rate and motion across subjects and at most a small within-subject increment that does not survive the ECG, and the EEG still contains cardiac signal after artefact correction. The simplest reading is physiological: heart-rate fluctuations drive BOLD changes and enter the EEG through residual pulse artefacts and autonomic covariation. EEG--BOLD coupling exists \citep{goldman2002simultaneous,laufs2003eeg}, but in these data a predictive model cannot separate it from physiology, and a surrogate test over the BOLD cannot either.

\paragraph{EEG-free references need to be used as nulls in EEG-to-fMRI synthesis.} We suggest: report EEG-free predictive distributions, including a stimulus-locked one, next to every model; claim EEG-specific information only above a subject-swap null; test whether the predicted uncertainty changes with the input; show with a positive control that the model, trained as reported, recovers a known input-coupled signal; test the EEG's increment over heart-rate and motion models; and back every claim of absence with an effect size.

\paragraph{Limitations.} The analyses use two motor-imagery datasets from one laboratory (17 and 9 subjects) with block designs, and tests with $n=17$ have limited power, although the equivalence tests bound what we could have missed. The positive control injects a linear, spatially fixed signal, and we did not retune every method's learning rate. Respiration was not recorded and may account for part of what we attribute to heart rate.